\section{Global form of the Standard Model gauge group}
\label{sec:global-form}
% Status: cited/open. Global-form discussion is source-supported and constrains future completions.

The local Standard Model gauge algebra admits several global forms.  The gauge group may be written as
\begin{equation}
  \GSM=\frac{SU(3)_C\times SU(2)_W\times U(1)_Y}{\Gamma},
  \qquad
  \Gamma\subset \mathbb Z_6,
  \label{eq:gsm-global}
\end{equation}
with \(\Gamma=1,\mathbb Z_2,\mathbb Z_3,\mathbb Z_6\) as the standard possibilities \cite{TongLineOperators2017,HsinGomis2024}.  Local perturbative fields see the Lie algebra.  Line operators, magnetic sectors, theta-angle periodicities, and topological probes see the quotient.

This is relevant to the DeVries program for two reasons.  First, any compactification, endpoint, or \(G_2\) completion has to specify the global group it realizes.  Second, a pole-spectrum determinant for the electroweak sector has to coexist with the same line-operator spectrum as the low-energy Standard Model.  Thus a proposed \(\mathcal O_J\) comes with a global-form compatibility condition:
\begin{equation}
  \mathcal O_J
  \quad\text{is defined in a theory with}\quad
  \GSM=(SU(3)\times SU(2)\times U(1))/\Gamma .
  \label{eq:global-form-compatibility}
\end{equation}

Tong's line-operator analysis classifies how the quotient changes the electric, magnetic, and dyonic line sectors \cite{TongLineOperators2017}.  The quotient restricts admissible electric line operators and changes the allowed magnetic line operators through Dirac quantization.  It also changes theta-angle periodicities.  After electroweak symmetry breaking, the line sectors descend to electromagnetic probes, so \(\Gamma\) leaves a low-energy topological signature.

Hsin and collaborators propose a transport-style probe of this global form through a fractional topological response coefficient \cite{HsinGomis2024}.  The important point for this manuscript is structural: the global form is, in principle, observable through nonlocal or topological response data.  It supplies a completion constraint for the DeVries program, because a string or compactification model generally determines more than a local Lie algebra.

Grand unification provides the comparison class.  SU(5) and Spin(10) naturally package Standard Model representations in a way that favors the \(\mathbb Z_6\) quotient, and the subgroup \(S(U(2)\times U(3))\) realizes the quotient structure inside SU(5) \cite{BaezHuerta2010,Britto2021}.  This places the GUT lineage beside the DeVries pole reading: high-scale structural placement is already known in unified models, while the DeVries construction assigns its clean value to the low-energy pole ratio and therefore needs a separate pole-placement derivation.

The global-form constraint can be expressed as a commutative diagram requirement.  If a compactification or endpoint model gives an ultraviolet group \(G_{\rm UV}\), then the low-energy electroweak sector should fit into
\begin{equation}
\begin{array}{ccc}
G_{\rm UV} & \longrightarrow & \GSM \\
\downarrow & & \downarrow \\
\mathcal O_{\rm UV} & \longrightarrow & \mathcal O_J ,
\end{array}
\label{eq:global-diagram}
\end{equation}
where the upper arrow fixes \(\Gamma\) and the lower arrow reduces the ultraviolet spectral operator to the DeVries two-channel operator.  A model that supplies \(\mathcal O_J\) while leaving \(\Gamma\) unspecified has an incomplete topological sector.

The mass-ratio argument itself should use local electroweak fields and pole positions.  The global-form material acts as a filter on completions.  A viable string/Kaluza--Klein/\(G_2\) derivation should state:
\begin{enumerate}
\item which quotient \(\Gamma\) appears;
\item which line operators survive compactification and electroweak breaking;
\item how the Higgs/order-parameter field transforms under the chosen global group;
\item how the DeVries operator respects the same quotient.
\end{enumerate}
